% ============================================================
% Chapter 5: Development
%   Section: Writing, Reviewing and Publishing an Article
% Included from chapters/05-development/development.tex
% Source: docs/decisions/writing.md (built 2026-10-15, merged as
% feature/newsroom-writing, outside the Oct 14 freeze by the author's
% decision).
% ============================================================

\section{Writing, Reviewing and Publishing an Article}
\label{sec:dev-writing}

The newsroom's last steps were built on 15~October~2026, outside the
feature freeze of 14~October by the author's decision; until then
writing from verified sources was listed as work for ``only if time is
left'' (Table~\ref{tab:priorities}). After an article has been analysed
and judged publishable, a person can ask for a new article on the same
story: the system finds news related to it, writes a draft from those
sources and the fact-checks in Catalan, Spanish or English following
the person's indications, reviews it, rewrites it if its metrics fall
short, and publishes it once a person approves it --- today only to the
system's own public reader.

\paragraph{One draft, one store.} Every stage reads and writes one
draft record, kept as a JSON file per draft beside the lake. Its body
cites its sources as \texttt{[n]}, numbered once and never reordered,
so a citation always points at the source it was written from. A
draft's editorial status is \emph{draft}, \emph{approved} or
\emph{published}; long steps run in the background, one at a time per
draft, and a person's edit sends an approved draft back to
\emph{draft}.

\paragraph{Writing, in a closed world.} Only an article the reader
publishes can be written from: analysed, admitted, and not judged
\texttt{FALSE} or \texttt{MISLEADING}. Related news comes from two
routes at once --- the internal corpus, by the analysed article's own
stored vector, and the open web, by its headline through the same
metasearch engine the fact-checker uses --- filtered to one source per
domain and none from the article's own outlet. The analysed article is
always source~\texttt{[1]}; then the related news; then only the
evidence a verdict was actually based on, each excerpt taken around
the passage the verdict quoted. One call to the language model writes
the draft, under rules that keep it in a closed world: only the
numbered sources and the fact-checks, a citation on every factual
sentence, no number, name, date or quote the sources lack, and the
verdicts deciding what may be stated --- a \texttt{FALSE} or
\texttt{UNVERIFIED} claim never as fact. Text in a source is material,
never an instruction. Checks without a model then repair what can be
repaired (citations to sources that do not exist are removed) and log
what needs a judgement: uncited sentences, runs copied from the
analysed article, numbers no source contains.

\paragraph{Reviewing against what the writer was given.} The review
scores nine metrics from 0 to 100: the title and the subtitle, and the
corrector's seven (Section~\ref{subsec:benchmark-writing}). Unlike free
text, a draft can be held to its sources, so the model is given each
source's excerpt and the verdicts and asked to check every factual
sentence against the source it cites; deterministic checks then bind
the scores where a small model reading a long prompt is unreliable ---
a factual sentence with no citation caps the hallucination score, a
citation to no source lowers factual consistency. A second call
proposes a correction, which is offered to the person only if it adds
no figure, keeps every citation on its statement, cites no new source,
stays in the draft's language and keeps a plausible length. Catalan is
supported with its own complete word lists and a named stand-in for
readability, since no Catalan reading-ease formula was found to cite.

\paragraph{Rewriting towards targets.} When metrics fall below their
targets --- 85 for factual consistency and hallucination, 80 for
grammar, lower for style and SEO, reasoned rather than fitted --- the
person can ask for a rewrite guided by the failing metrics. Each round
is a closed-world call over the draft's own sources, the same checks as
the correction, and a new review; a rewrite that breaks a check is not
kept, and any earlier text can be restored.

\paragraph{Publishing, with the AI disclosed.} An approved draft is
published as a snapshot built from an allowlist of fields --- never the
instructions, the review or any source's excerpt, which is somebody
else's text --- so the public page shows what a person approved until
they publish again or take it down. Every published article says,
before its text, that it was written by an AI model, names the model,
and says that a person approved it. The EU AI Act's Article~50(4),
applicable from 2~August~2026, asks for such a disclosure for
AI-generated text published to inform the public, with an exemption for
text under human editorial responsibility; whether this site meets the
exemption has not been judged, and disclosing does not depend on it.
Social channels are listed with what each would need (both publish
media, not text) and are not built.

\paragraph{What was measured.} The writer was run directly, outside the
stack, with \texttt{llama3.2:3b} on four drafts of a sample article in
Catalan and Spanish: none of the four cited a source; three stated as
fact a claim the fact-checks had left \texttt{UNVERIFIED} (the planted
\texttt{FALSE} one was not stated); and the Catalan had errors a reader
sees at once. A fifth run repeated it. The review, on two seeded drafts,
caught a claim judged \texttt{FALSE} and the missing subtitle, but
missed both spelling errors and a wrong figure, and its correction
changed a fact and moved a citation --- which is how the figure and
moved-citation checks came to exist. The rewrite checks stopped four of
five rewrites, three of them for pinning a source onto the false claim;
they cannot see a fact a rewrite leaves out. Two practical findings
followed: the 3-billion-parameter model is fit to verify claims quickly
on a CPU and not to write or correct an article, and Ollama silently
cuts a long prompt to its default 4,096-token context, so the writer's
model needs its own larger context. On 19~October two openly licensed
writers were tried the same way: Qwen3.5 at 4 and 9 billion parameters
(Apache~2.0), served at a 16k-token context and with their reasoning
mode switched off, since both otherwise think before every answer. Both
cited a source for every sentence or all but one, filled the metadata
and stated neither the unverified nor the planted false claim; on a
processor alone a draft took about three minutes with the smaller and
four to six with the larger. Their Catalan differed: the smaller made
slips a reader notices, the larger wrote ``litigi'' (a lawsuit) for
\emph{liti} (lithium) throughout. The smaller is the recommended writer
(Subsection~\ref{subsec:results-writing-trial}); no draft has yet been
written through the running stack.
